\section{等势点与电路的拓扑}
\subsection{导线连接的点先合并}

理想导线的电阻为零，只要导线中电流有限，两端电势差就为零。复杂图形中的第一步往往不是计算，而是沿理想导线追踪节点，把同一节点在图中所有位置标为同一个字母。经过重新绘图，本来看似依次排列的电阻可能连接在同一对端点之间，实际属于并联。

重画电路可以拉长导线、移动电阻的位置，但不能改变每个元件两端的节点，也不能把几何交点自动变成电连接点。三维框架图投影到纸面上尤其要留意：投影中的交叉不等于空间中相接。每次重画后，可逐一核对“元件—起点—终点”，确认边数、接点和端口均未改变。

\fig{c3:topology-redraw}{追踪理想导线后，左侧三个电阻实际都跨接 $A,B$，重画后成为并联}{\begin{tikzpicture}

\begin{scope}[scale=.85]
\draw(0,0)to[R,l={$R_1$}](2,0)to[R,l={$R_2$}](4,0)to[R,l={$R_3$}](6,0);
\draw(0,0)--(0,1.2)--(4,1.2)--(4,0);\draw(2,0)--(2,-1.2)--(6,-1.2)--(6,0);
\foreach \x/\lab in {0/A,2/B,4/A,6/B}{\fill(\x,0)circle(1.5pt)node[below left]{$\lab$};}
\end{scope}
\begin{scope}[xshift=7cm,scale=.85]
\foreach \y/\lab in {1.2/R_1,0/R_2,-1.2/R_3}{\draw(0,\y)to[R,l={$\lab$}](3,\y);}
\draw(0,-1.2)--(0,1.2) (3,-1.2)--(3,1.2);
\node[left]at(0,0){$A$};\node[right]at(3,0){$B$};
\end{scope}

\end{tikzpicture}}
图\ref{c3:topology-redraw}中，沿左图上方导线到达的两点都标为 $A$，沿下方导线到达的两点都标为 $B$。$R_2$ 虽画在中间，仍连接这同一对节点；其画法方向与另外两只相反，不影响并联关系。重画后逐项核对 $R_1,R_2,R_3$ 的两个端点，能直接确认拓扑未变。

\subsection{几何对称为何能够推出等势}

给纯电阻网络两端固定电势后，内部节点电势由线性守恒方程决定。若一种置换或空间变换使端点、每条连线及其电阻值都保持不变，变换后的电势分布仍满足同一边界条件和方程。解的唯一性使原分布与变换后的分布相同，所以变换相互对应的节点等势。

唯一性可以从能量角度说明。假设同一边界条件存在两组解，取各节点电势之差 $w_i$，端点的差为零。对差分电流的节点方程乘以 $w_i$ 后求和。同一条边在两端节点各出现一次，两项合起来为
\[
 w_i\frac{w_i-w_j}{R_{ij}}+w_j\frac{w_j-w_i}{R_{ij}}
 =\frac{w_i^2-2w_iw_j+w_j^2}{R_{ij}}
 =\frac{(w_i-w_j)^2}{R_{ij}}.
\]
端口处 $w_i=0$，其外部电流不产生额外项，因而全部边相加给出
\[
 \sum_{\langle i,j\rangle}\frac{(w_i-w_j)^2}{R_{ij}}=0.
\]
所有电阻为正，各项非负，故连通边两端的 $w$ 都相等；端点为零又使整个连通网络的 $w$ 为零。因此解唯一。对称等势并不是凭外观作猜测，而是边界不变和方程唯一性的结果。

如果变换交换高、低电势端点而非分别保持，应把电势同时改成 $U-V$，才能得到同一边界。这时对应点通常满足 $V_P+V_{P'}=U$，未必等势。只有处在这种交换变换不动位置的节点，才可推出电势为 $U/2$。这个区分能避免在方形或立方体网络中误合并不同节点。


\fig{c3:swap-boundaries}{端点交换时应同时反演电势；对应节点的电势和为 $U$，一般并不等势}{\begin{tikzpicture}

\draw(0,0)to[R,l={$r$}](2,0)to[R,l={$r$}](4,0)to[R,l={$r$}](6,0)to[R,l={$r$}](8,0);
\foreach \x/\lab/\pot in {0/A/U,2/P/{3U/4},4/C/{U/2},6/Q/{U/4},8/B/0}{\fill(\x,0)circle(1.5pt)node[below]{$\lab$};\node[above]at(\x,.9){$\pot$};}
\draw[guide](4,-.6)--(4,1.4);\node[below]at(4,-1){反射交换 $A,B$ 与 $P,Q$，$C$ 不动};

\end{tikzpicture}}
图\ref{c3:swap-boundaries}的反射把高、低电势端互换，故 $V_P+V_Q=U$；$P,Q$ 的电势分别为 $3U/4,U/4$，不能合并。中点 $C$ 在反射下不动，才满足 $V_C=U-V_C$，从而 $V_C=U/2$。判断对称时必须把端口边界一并变换。

\section{等势合并、支路断开与分裂}
\subsection{允许的简化应由电势条件支持}

两节点在当前端口激励下严格等势时，用理想导线连接二者不会引入额外电流，因此不改变该端口的电压电流关系。若它们之间原有有限电阻，原电流就是零，删除该电阻也不改变结果。这两种操作分别是等势合并与无流支路断开。

反过来，把一个节点拆成两个节点并断开，需要证明断开后两节点仍由各自网络约束为同电势。仅仅因为断开前二者属于同一节点，并不能保证断开后仍然等势。常见可行情况是两条相似分压支路的中点，断开后仍有相同分压比；若两侧电阻比例不同，就会出现电势差，断开操作改变原电路。

等势简化也与当前端口有关。立方体电阻网络在体对角线端点加电压时，可按距离高电势端的层数合并节点；改从相邻顶点供电后，同一批节点一般不再等势。解答应明确谁是输入、输出端以及哪些对称操作保持它们，而不能把某张简化图当作立方体网络的永久结构。

\fig{c3:equal-potential-merge}{反射保持 $A,B$ 不动且互换 $M,N$；对称条件成立时中间支路无流}{\begin{tikzpicture}

\draw(0,0)to[R,l={$R_1$}](2.5,1.7)to[R,l={$R_2$}](5,0)to[R,l_={$R_2$}](2.5,-1.7)to[R,l_={$R_1$}](0,0);
\draw(2.5,1.7)to[R,l={$R_0$}](2.5,-1.7);
\foreach \x/\y/\lab in {0/0/A,5/0/B,2.5/1.7/M,2.5/-1.7/N}{\fill(\x,\y)circle(1.5pt)node[above right]{$\lab$};}
\node[right]at(5.7,1){$V_M=V_N$};
\node[right]at(5.7,0){$I_{MN}=0$};
\node[right]at(5.7,-1){可删 $R_0$ 或合并 $M,N$};

\end{tikzpicture}}
图\ref{c3:equal-potential-merge}中，上下两臂的连接和阻值逐项对应。删去 $R_0$ 后，两个中点仍分别由相同分压比决定，等势条件保持；若改动其中一臂，便须先重新求两中点电势，原有删支路或分裂操作不再自动成立。

\subsection{等势层中的并联数目}

若网络被分为若干等势层，相邻两层之间共有 $m$ 条相同电阻 $r$，它们跨接同一对电势，所以并联等效为 $r/m$。不同层之间通过同一总电流，层间等效电阻再串联。最重要的计数是跨越两层的边数，而不是层内节点数。

\fig{c3:cube-unreduced}{体对角线端口的原立方体；每棱为 $r$，未标点的交叉处不连接，二进制数中 $1$ 的个数决定分层}{\begin{tikzpicture}

\coordinate(A)at(0,0);\coordinate(P)at(3,0);\coordinate(Q)at(1,1);\coordinate(R)at(0,3);
\coordinate(S)at(4,1);\coordinate(T)at(3,3);\coordinate(W)at(1,4);\coordinate(B)at(4,4);
\foreach \u/\v in {A/P,A/Q,A/R,P/S,P/T,Q/S,Q/W,R/T,R/W,S/B,T/B,W/B}{\draw(\u)to[R](\v);}
\fill[red](A)circle(2pt)node[below left]{$A:000$};\fill[red](B)circle(2pt)node[above right]{$B:111$};
\foreach \n/\lab in {P/100,Q/010,R/001}{\fill[blue](\n)circle(2pt)node[below left]{$\lab$};}
\foreach \n/\lab in {S/110,T/101,W/011}{\fill[orange!80!black](\n)circle(2pt)node[above right]{$\lab$};}

\end{tikzpicture}}
图\ref{c3:cube-unreduced}中，坐标置换保持 $000,111$ 两端不动，三个蓝点互相对应，三个橙点也互相对应。$A$ 到蓝层有三条边；每个蓝点分别通向两个橙点，中间共六条边；橙层到 $B$ 又有三条边。等势分组来自网络对称，六条跨层边则来自实际连接的计数。


以立方体体对角线为端点，四层节点数为 $1,3,3,1$，相邻层的边数为 $3,6,3$，所以图\ref{fig:c3:layers}的串联简化给出 $r/3+r/6+r/3=5r/6$。中间两层各有三个节点，却有六条连接边；将其误计为三条会把中段电阻算大一倍。

\typeout{ECFIG:fig:c3:layers}\fig{fig:c3:layers}{立方体体对角线供电的等势层}{\begin{tikzpicture}[scale=.88]
\draw(0,0)to[R,l={$r/3$}](2,0);\draw(2,0)to[R,l={$r/6$}](4,0);\draw(4,0)to[R,l={$r/3$}](6,0);\node[left]at(0,0){$A$};\node[right]at(6,0){$B$};\node at(1,-1){3 条};\node at(3,-1){6 条};\node at(5,-1){3 条};
\end{tikzpicture}}

\section{完全图与高维网络}
\subsection{完全连接的节点群}

在 $N$ 个节点两两以电阻 $r$ 连接的网络中，取其中两个节点 $A,B$ 为端点，其余 $N-2$ 个节点在交换彼此时网络和端口都不变，故它们等势。交换 $A,B$ 后它们的整体位置不变，又使电势变为 $U-V$，所以每个内部节点电势均为 $U/2$。

内部节点之间没有电流。$A,B$ 的直接支路为 $r$，每个内部节点提供一个 $2r$ 的间接通道，这些通道并联，因而
\[
 \frac1{R_{AB}}=\frac1r+\frac{N-2}{2r}=\frac N{2r},\qquad
 R_{AB}=\frac{2r}{N}.
\]
几何上可以把节点画成正多面体，也可画在圆周上；决定结果的是连接关系和阻值，不是线段的实际长度。若题设规定电阻与几何长度成正比，就必须另外保留每条边的真实阻值，不能继续使用“所有边为 $r$”的公式。

\fig{c3:complete-graph-reduction}{内部节点等势后，它们之间的边无电流；保留直接通道和经内部节点的间接通道}{\begin{tikzpicture}

\draw(0,0)to[R,l={$r$}](6,0);
\draw(0,0)--(0,2)to[R,l={$r$}](3,2)to[R,l={$r$}](6,2)--(6,0);
\draw(0,0)--(0,-2)to[R,l_={$r$}](3,-2)to[R,l_={$r$}](6,-2)--(6,0);
\node[left]at(0,0){$A:U$};\node[right]at(6,0){$B:0$};
\fill(3,2)circle(1.5pt)node[above]{$U/2$};\fill(3,-2)circle(1.5pt)node[below]{$U/2$};
\node[right]at(6.5,-1.5){间接通道共 $N-2$ 条};\node at(3,-.8){$\cdots$};

\end{tikzpicture}}
图\ref{c3:complete-graph-reduction}的每条间接通道含两个 $r$ 串联，且与直接的 $r$ 同跨 $A,B$。原内部节点间的电阻已因零电流省去，它们对本端口的电流没有贡献。由通道电导相加得到 $N/(2r)$，也说明增加完全连接节点数会降低该端口电阻。

\subsection{二进制标记与层间计数}

$d$ 维超立方体的节点可用 $d$ 位二进制数标记，两节点仅一位不同就有一条电阻 $r$ 相连。以全零和全一为端点，对坐标作任意置换不改变网络，因此含相同数量的 $1$ 的节点等势。

第 $k$ 层有 $\binom dk$ 个节点，每个节点有 $d-k$ 条边通向第 $k+1$ 层，故跨层边数是 $(d-k)\binom dk$。不同节点出发的边不会重复计数，因为每条跨层边都只有一个第 $k$ 层端点。于是
\[
 R_d=r\sum_{k=0}^{d-1}\frac1{(d-k)\binom dk}.
\]
该式是体对角线端口的结果。它把三维网络中的“等势层—并联边—串联层”方法推广到更高维，也提供了清晰的组合计数练习。求其他端口时，应重新确定保持端口的对称群，不能只改变层数的上限。


可用低维结果检查超立方体公式。$d=1$ 时只有一条边，$R_1=r$；$d=2$ 时两条 $2r$ 路径并联，$R_2=r$；$d=3$ 时跨层边数是 $3,6,3$，回到图\ref{c3:cube-unreduced}的 $5r/6$。例如 $d=4$ 时跨层边数为 $4,12,12,4$，故 $R_4=2r/3$。这些系数来自边的数目，不能只用每层节点数代替。

\section{星形与三角形的端口等效}
\subsection{三个端口需要三个独立参数}

一个两端无源电阻网络只需一个等效电阻描述，三个端口的网络则必须保留三对端口的响应。把三个端口的网络直接替换成单个电阻，会丢失第三端接入后对电流分配的影响。星形和三角形都是用三个电阻表示三端网络，等效时应在第三端悬空条件下，使任意两端的电阻分别相同。

\fig{fig:c3:delta}{三角形的三个端口}{\begin{tikzpicture}[scale=.88]
\draw(0,0)to[R](4.6,0)to[R](2.3,4)to[R](0,0);\node at(2.3,-.5){$R_{12}$};\node at(4.2,2.1){$R_{23}$};\node at(.35,2.1){$R_{31}$};\node[left]at(0,0){$1$};\node[right]at(4.6,0){$2$};\node[above]at(2.3,4){$3$};
\end{tikzpicture}}

设三角形边为 $R_{12},R_{23},R_{31}$，总和为 $S$，星形三臂为 $r_1,r_2,r_3$。\fig{fig:c3:star}{保留相同端口的星形}{\begin{tikzpicture}[scale=.88]
\draw(2.3,1.6)to[R](0,0);\draw(2.3,1.6)to[R](4.6,0);\draw(2.3,1.6)to[R](2.3,4.2);\node at(.75,1.35){$r_1$};\node at(3.85,1.35){$r_2$};\node at(1.7,3){$r_3$};\node[left]at(0,0){$1$};\node[right]at(4.6,0){$2$};\node[above]at(2.3,4.2){$3$};\node[right]at(2.3,1.6){$O$};
\end{tikzpicture}}
图\ref{fig:c3:delta}与图\ref{fig:c3:star}保留同样的 $1,2,3$ 三端。比较端口 $1,2$ 时，三角形的直接边 $R_{12}$ 与串联路径 $1\to3\to2$ 并联；星形的第三端悬空，使 $r_3$ 无电流，只剩 $r_1,r_2$ 串联。


第三端悬空时，
\[
 r_1+r_2=\frac{R_{12}(R_{23}+R_{31})}{S},\quad
 r_2+r_3=\frac{R_{23}(R_{31}+R_{12})}{S},
\]
\[
 r_3+r_1=\frac{R_{31}(R_{12}+R_{23})}{S}.
\]
第一式与第三式相加再减第二式，左边为 $2r_1$；右边合并后为 $2R_{12}R_{31}/S$，所以
\[
 r_1=\frac{R_{12}R_{31}}{S},\quad
 r_2=\frac{R_{12}R_{23}}{S},\quad
 r_3=\frac{R_{23}R_{31}}{S}.
\]
每个星臂对应的是原三角形中与该端口相邻的两边之积，分母是三边之和。画图标清端口比背诵下标更可靠。

\typeout{ECFIG:fig:c3:delta}
\typeout{ECFIG:fig:c3:star}
\subsection{反变换与使用位置}

由上述关系可反解三角形边。例如
\[
 \frac{r_1r_2+r_2r_3+r_3r_1}{r_3}
 =r_1+r_2+\frac{r_1r_2}{r_3}=R_{12}.
\]
把三角形到星形的表达式代入，$r_1r_2/r_3=R_{12}^2/S$，而 $r_1+r_2=R_{12}(R_{31}+R_{23})/S$，两项之和正好为 $R_{12}$。所以反变换并非凭相同的几何外形猜比例，而是由同一组端口测量关系反解。

循环交换下标即得到另两边。三条相等三角形电阻 $R$ 对应三条 $R/3$ 星臂；三条相等星臂 $r$ 对应三条 $3r$ 三角形边。两种方向的比例相反，计算前应确定是在消去中心节点，还是在引入中心节点。

选择变换位置时，应关注变换后能否立即形成串并联，而不是把所有三角形都逐个替换。星三角变换只保留端口外部响应，变换后新增的中心节点不是原网络中的物理节点；原某电阻上的电流也不能直接等同于变换后某星臂电流。若题目还要求内部电流，应用变换求得端口电势后回到真实电路计算。


\section{复杂网络的核验与方法选择}
\subsection{从整体界限检查结果}

纯电阻网络的等效电阻必须为正。向网络添加一条正电阻支路不可能使等效电阻增大，删除支路不可能使其减小。下面用固定端电压 $U$ 时的耗散功率解释这一结论。记真实节点电势为 $V_i$，另取一组满足同一端口电势的试探值 $V_i+w_i$，其中端口处 $w_i=0$。将每条边的功率展开，
\[
 P[V+w]-P[V]
 =2\sum_{\langle i,j\rangle}\frac{(V_i-V_j)(w_i-w_j)}{R_{ij}}
  +\sum_{\langle i,j\rangle}\frac{(w_i-w_j)^2}{R_{ij}}.
\]
第一项按节点收集后，是 $2\sum_i w_i\sum_j(V_i-V_j)/R_{ij}$。内部节点的电流和为零，端口的 $w_i$ 为零，故这一项消失。剩下的是非负平方和，说明真实电势使固定边界下的总功率最小。

若在节点 $p,q$ 之间添加电阻 $R_0>0$，任意同一组试探电势下均有
\[
 P_{\text{新}}[V]=P_{\text{原}}[V]+\frac{(V_p-V_q)^2}{R_0}
 \ge P_{\text{原}}[V].
\]
因此最小值也满足 $P_{\text{新},\min}\ge P_{\text{原},\min}$。端口功率为 $UI=U^2/R_{\text{等效}}$，所以 $R_{\text{新}}\le R_{\text{原}}$。若新增支路原本就连接两个等势点，例如图\ref{c3:equal-potential-merge}中的 $M,N$，新增功率项为零，原电势继续满足全部节点方程，电阻恰好不变。这也给出单调性取等号的具体情形。

对称法给出简洁结果时，仍应选一组节点电势或层间电流回查守恒。对于不能完全合并的网络，保留少量未知节点继续用基尔霍夫方程即可；不必强行把每道题化为纯串并联。图形重画、对称合并、端口等效和星三角变换是相互衔接的方法，选用它们的共同目标是减少未知量，同时保留全部端口约束。


\section{虚拟注入电流与有限网络}
\subsection{构造便于对称分析的辅助电流}

一组假想电流可以在若干节点流入、另一些节点流出，只要流入和流出的总代数和为零，就能作为同一线性网络上的辅助问题。求出两个辅助问题后，电流与电势可相加；若不需要的节点注入恰好抵消，就得到真正的两端供电问题。这种构造不是向真实题目增加元件，而是在同一组电阻上分解外部电流边界。

对具有 $N$ 个等价节点的网络，先在 $A$ 注入 $I$，从其余每个节点流出 $I/(N-1)$。再作另一辅助问题：在 $B$ 流出 $I$，从其余每个节点注入 $I/(N-1)$。叠加后除了 $A,B$，其余节点外部电流都抵消，真实端口电流为
\[
 I_{\text{端}}=I+\frac{I}{N-1}=\frac{N}{N-1}I.
\]
如果网络对称性使两个辅助问题的 $A,B$ 电势差相同，真实电势差就是其中一个的两倍。求等效电阻时必须除以叠加后的 $I_{\text{端}}$，不能仍除以原辅助电流 $I$。

超立方体适合这种方法。全零点注入电流后，同一汉明距离层节点等势；第 $k$ 层到下一层的总电流，等于总注入电流减去前面已经流出的各节点电流。每层有多少外部抽取、跨层共有多少电阻，都应在过程图中标清。与两端直接供电相比，辅助问题恢复了围绕注入点的坐标置换对称性。


\fig{c3:virtual-injection}{图示同一个四节点环的两辅助边界叠加；电阻不变，只组合外部注入电流}{\begin{tikzpicture}

\foreach \shift/\mode in {0/1,4.4/2,8.8/3}{\begin{scope}[xshift=\shift cm,scale=.7]
\draw(0,0)to[R](2,0)to[R](2,2)to[R](0,2)to[R](0,0);
\node[above]at(0,2){$A$};\node[above]at(2,2){$B$};
\node[left]at(0,0){$P$};\node[right]at(2,0){$Q$};
\ifnum\mode=1
\draw[flow](-1,2)--(0,2)node[pos=0,left]{$I$};\draw[flow](2,2)--(3,2)node[right]{$I/3$};
\draw[flow](0,0)--(0,-1)node[below]{$I/3$};\draw[flow](2,0)--(2,-1)node[below]{$I/3$};
\fi
\ifnum\mode=2
\draw[flow](-1,2)--(0,2)node[pos=0,left]{$I/3$};\draw[flow](2,2)--(3,2)node[right]{$I$};
\draw[flow](0,-1)--(0,0)node[pos=0,below]{$I/3$};\draw[flow](2,-1)--(2,0)node[pos=0,below]{$I/3$};
\fi
\ifnum\mode=3
\draw[flow](-1,2)--(0,2)node[pos=0,left]{$4I/3$};\draw[flow](2,2)--(3,2)node[right]{$4I/3$};
\node[below]at(1,-.9){$P,Q$ 外部电流抵消};
\fi
\end{scope}}

\end{tikzpicture}}
图\ref{c3:virtual-injection}取 $N=4$，前两图叠加后，$P,Q$ 不再有外部抽取或注入，$A,B$ 的端口电流却增为 $4I/3$。求电阻必须用合成后的电势差除以这个总电流。各辅助图中的电流不是可以忽略的“额外支路”，而是线性方程右端的不同边界条件。

\section{多端网络等效的限制}
\subsection{四端网络不能普遍只用一个星形表示}

三端无源互易网络可用三个端间测量电阻确定星形或三角形参数。但四个端点有六种端口对，普遍的四端网络需要更多独立参数。仅有四条臂的中心星形只有四个参数，不能代表任意六参数的四端网络。若某一特殊对称构型可以换成星形，必须核对全部需要的端口关系，不能只核对其中一对就推广。

对带多个外接端点的子网络，内部替换要保持任意外部激励下所有端口电压与注入电流的关系。只在一个指定端口测得相同的等效电阻，不保证它嵌入另一个电路后仍可任意替代。这也是无限嵌套正方形不能直接把内图当作二端电阻、丢弃其余两个连接点的原因。

\fig{c3:four-terminal-obstruction}{左侧四端环一般不能换成右侧单中心四臂星形；四端等效须保留所有端口关系}{\begin{tikzpicture}

\begin{scope}
\draw(0,0)to[R,l={$r$}](3,0)to[R,l={$r$}](3,3)to[R,l={$r$}](0,3)to[R,l={$r$}](0,0);
\foreach \x/\y/\lab in {0/0/1,3/0/2,3/3/3,0/3/4}{\fill(\x,\y)circle(1.5pt)node[above right]{$\lab$};}
\end{scope}
\begin{scope}[xshift=6cm]
\foreach \x/\y/\lab in {0/0/1,3/0/2,3/3/3,0/3/4}{\draw(1.5,1.5)to[R,l={$r_{\lab}$}](\x,\y);\node[above right]at(\x,\y){$\lab$};}
\fill(1.5,1.5)circle(1.5pt);
\end{scope}

\end{tikzpicture}}
图\ref{c3:four-terminal-obstruction}右侧若其余端悬空，任意两端测得 $R_{ij}=r_i+r_j$，因而必须有
\[
 R_{12}+R_{34}=R_{13}+R_{24}=R_{14}+R_{23}.
\]
左侧四边均为 $r$ 时，相邻两端电阻为 $3r/4$，相对两端电阻为 $r$。于是第一组和为 $3r/2$，第二组为 $2r$，不满足上述必要条件，不能用四条星臂普遍替换。

\subsection{边中点作为端口时要保留半段电阻}

端口落在一条均匀电阻导线的中点时，必须先把该棱拆为两个 $r/2$，再按真正节点求解。如果两个端口之间的棱段被理想导线连接，应先合并；若只是投影相交则不能连接。对称分组必须同时保留端口位置和半段阻值，否则沿用顶点端口的结果会少算串联部分。


\fig{c3:mid-edge-port}{均匀边上的中点端口必须成为新节点；两半边各为 $r/2$}{\begin{tikzpicture}

\draw(0,0)to[R,l={$r$}](4,0);\node[left]at(0,0){$A$};\node[right]at(4,0){$B$};
\node at(5.3,0){$\Longrightarrow$};
\begin{scope}[xshift=6.6cm]
\draw(0,0)to[R,l={$r/2$}](2,0)to[R,l={$r/2$}](4,0);
\node[left]at(0,0){$A$};\node[right]at(4,0){$B$};\fill(2,0)circle(2pt)node[below]{$P$};
\draw[flow](2,-1.2)--(2,-.2)node[pos=0,below]{端口电流};
\end{scope}

\end{tikzpicture}}
图\ref{c3:mid-edge-port}中，外部电流可在新节点 $P$ 流入或流出，两半边不再必然通过相同电流。因此应先拆边、标记真实端口，再寻找保持这些端口的对称变换，不能沿用整条棱仅作为内部元件时的简化。
